% !PRINTED 143
\weberpage{143}

Man kann das Formelsystem (8) auch nach der entgegengesetzten Richtung fortsetzen, wenn man die Gleichungen
\[
 S[\alpha^{-1}f(\alpha)]=0,\qquad
 S[\alpha^{-2}f(\alpha)]=0,\ldots
\]
bildet. Man erhält dadurch ein Mittel, um die Potenzsummen $s_\nu$ auch für negative Exponenten $\nu$ durch die Grundfunctionen auszudrücken. Diese Summen der negativen Potenzen gehören gleichfalls zu den symmetrischen Functionen, wenn auch nicht mehr zu den \gesperrt{ganzen}, sondern zu den \gesperrt{gebrochenen}. Sie gehen erst durch Multiplication mit Potenzen des Productes $\alpha_1\alpha_2\cdots\alpha_n$ in ganze Functionen über. Der Vollständigkeit wegen setzen wir die zwei ersten dieser Formeln hierher:
\begin{equation}
\begin{aligned}
0&=s_{n-1}+a_1s_{n-2}+a_2s_{n-3}+\cdots+na_{n-1}+a_ns_{-1},\\
0&=s_{n-2}+a_1s_{n-3}+a_2s_{n-4}+\cdots+a_{n-1}s_{-1}+a_ns_{-2}.
\end{aligned}
\tag{11}
\end{equation}
die sich nach (7) auch so darstellen lassen:
\[
 a_{n-1}+a_ns_{-1}=0,\qquad 2a_{n-2}+a_{n-1}s_{-1}+a_ns_{-2}=0.
\]

\begin{center}
{\large \S.\ 43.}

\vspace{0.5em}
{\large\gesperrt{Beweis des Hauptsatzes für zwei Variable.}}
\end{center}

Wir gehen nunmehr zum Beweis des Fundamentalsatzes der Theorie der symmetrischen Functionen über, dass sie alle rational durch die symmetrischen Grundfunctionen ausdrückbar sind. Da wir die vollständige Induction als Beweismittel anwenden, so leiten wir den Satz zunächst unter der Voraussetzung ab, dass nur zwei unabhängige Veränderliche $\alpha,\beta$ gegeben seien, aber auf einem Wege, der zugleich für den allgemeinen Beweis den leitenden Gedanken hervortreten lassen wird.

Wir bezeichnen die symmetrischen Grundfunctionen mit
\begin{equation}
 a=-(\alpha+\beta),\qquad b=\alpha\beta
\tag{1}
\end{equation}
und setzen demgemäss
\begin{equation}
 f(x)=(x-\alpha)(x-\beta)=x^2+ax+b.
\tag{2}
\end{equation}

Es sei nun $S(\alpha,\beta)$ irgend eine ganze rationale und symmetrische Function von $\alpha$ und $\beta$. Wir können für $\beta$ aus (1) den Werth $-(a+\alpha)$ einsetzen und erhalten, wenn wir nach Potenzen von $\alpha$ ordnen,
\begin{equation}
 S(\alpha,\beta)=S(\alpha,-a-\alpha)=A_0\alpha^m+A_1\alpha^{m-1}+\cdots+A_{m-1}\alpha+A_m,
\tag{3}
\end{equation}
%CONT>>
